% SPDX-License-Identifier: Apache-2.0
\section{Evaluation}
\label{sec:evaluation}

The preceding sections argue that three kinds of evidence --- machine-checked
proofs, schema-validated bills of materials, and cross-implementation
agreement --- each answer a question narrower than the one a reader takes them
to answer. This section measures that claim on one project rather than
illustrating it: how many declarations the checker could actually resolve, what
each class of check can and cannot see when the same artefact is disturbed, and
what happened when the same failure mode was looked for one layer below the
document, in the delivery path that puts evidence on the network.

All figures were produced by the commands named in the reproduction appendix
(\texttt{verifact verify} for the counts, \texttt{run-e2e.cjs} for the
disturbance matrix, \texttt{check-live.cjs} for the delivery gate), each at
the commit recorded in the provenance note to this section; none are quoted
from an earlier revision of this paper. The corpus counts here are measured on
clean checkouts at \texttt{b2910d5} (product repository) and \texttt{163efe3}
(\texttt{verifact}); the CBOM figures in Section~\ref{sec:case-study-cbom} are
measured at \texttt{95056b4}.

\subsection{What the checker could resolve}

The public corpus holds 2{,}337 statements that assert a property of a named
artefact (a file, symbol, constant or derived parameter). Of these, 2{,}127 are
claims about a code artefact the checker can locate. The other 210 fall outside
that scope for five different reasons, and the table separates them rather than
lumping them: standard references that are not artefact claims (25), a sealed
evidence directory that must not be edited (84), claims ignored by configuration
(21), claims dispositioned as narrative (72), and claims denied by policy (8).
Of the 2{,}127, a binding could be attempted for 2{,}109: 136 verified, 1{,}759
unbound and 214 ambiguous. The remaining 18 anchor to a standard reference
rather than to an artefact, so they are locatable but not bindable.

\begin{table}[t]
\centering
\small
\begin{tabular}{@{}lr@{}}
\toprule
Measure & Count \\
\midrule
Declarations in the public corpus                     & 2{,}337 \\
\quad claims about a locatable code artefact          & 2{,}127 \\
\qquad verified (binding found and consistent)        &   136 \\
\qquad unbound (no binding found)                     & 1{,}759 \\
\qquad ambiguous (multiple candidate bindings)        &   214 \\
\qquad anchor is a standard reference, not an artefact &    18 \\
\quad claims outside artefact scope                   &   210 \\
\qquad standard references, not artefact claims       &    25 \\
\qquad sealed evidence directory                      &    84 \\
\qquad ignored by configuration                       &    21 \\
\qquad dispositioned as narrative                     &    72 \\
\qquad denied by policy                               &     8 \\
\midrule
Drift between declaration and artefact               &     0 \\
Explicit exemptions, with author and expiry          &    19 \\
\bottomrule
\end{tabular}
\caption{Corpus and verdict distribution. Every row closes against the rows
beneath it: $136 + 1{,}759 + 214 + 18 = 2{,}127$, and $2{,}127 + 210 = 2{,}337$.
The tool's own \texttt{allJudgedClaims} is 2{,}207, which is $2{,}127 + 72 + 8$;
it is not the number of verdicts, because 80 of those claims are dispositioned or
denied rather than judged, so it is not used as a denominator here.}
\label{tab:verdicts}
\end{table}

Three denominators can be defended, and the difference between them is smaller
than it is usually made to sound. $136/2{,}127$ is the rate over claims the
checker can locate; $136/2{,}109$ is the rate over claims whose binding it could
attempt. Both are $6.4\%$. Quoting either alone would hide the 1{,}759, which is
the number that says what the method does not cover; quoting the tool's
$2{,}207$ would be worse, because it mixes judged and unjudged claims.

The interesting number in Table~\ref{tab:verdicts} is not 136 but 1{,}759. A
checker that reported only ''136 verified'' would be read as a statement about
the corpus; what it actually establishes is that for 1{,}759 declarations the
corpus does not record enough to decide. Those are not failures of the
declarations. Most are prose assertions whose binding was never written down.
They are the boundary of what the method covers, and they are invisible in a
green summary line.

The 19 exemptions were re-tested by removing all of them and re-running the
checker: 6 then appeared as drift, and 13 did not, because they describe a
different measurement basis or a property that is not an artefact attribute.
Reporting ''19 exemptions'' without that split would have overstated the residue
by a factor of three.

\subsection{What each class of evidence can see}
\label{sec:eval-baseline}

The same claim, the same artefact and the same four disturbances were presented
to two checkers: a value-level checker that compares the numbers a document
states against the numbers the artefact yields, and a byte-level checker that
compares the artefact against a digest list sealed by a timestamp token.

\begin{table}[t]
\centering
\small
\begin{tabular}{@{}clll@{}}
\toprule
 & Disturbance & Value-level checker & Byte-level checker \\
\midrule
S0 & none                              & pass, exit 0        & pass, exit 0 \\
S1 & one number changed (62 $\to$ 61)  & drift, exit 1       & drift, exit 1 \\
S2 & one byte flipped, no number       & \textbf{pass, exit 0} & drift, exit 1 \\
S3 & S2 + digest list rewritten        & pass, exit 0        & drift, exit 1 (imprint) \\
S4 & signing CA omitted                & n/a                 & unverifiable, exit 3 \\
\bottomrule
\end{tabular}
\caption{Two checkers, four disturbances. S2 is the row that motivates this
paper: the value-level checker does not warn, it reports success, because the
document still says what it said. S3 shows why the digest list cannot be its own
authority.}
\label{tab:s0s4}
\end{table}

Row S2 is not a corner case. Editing bytes without editing statements is what a
line-ending conversion, a re-encoding, a minifier or a packaging step does by
default, and every one of those steps is routine. Row S3 is the reason the
digest list is sealed rather than trusted: rewriting the list restores link one
of the chain, and only the token's imprint over the original list catches it.
Row S4 exists to keep ''could not decide'' from rendering as ''passed''. The
byte-level checker returns a distinct exit code, not a warning on a success
path.

\subsection{The deploy that returned zero}
\label{sec:eval-deploy}

The three contributions above concern evidence that is produced and read. The
same failure mode occurs one layer lower, in the path that publishes the
evidence, and it is included here because it was measured in our own
infrastructure rather than hypothesised.

The deploy flow in use ended with \texttt{git reset -{}-hard origin/main} followed
by a local health check on the application port. Both returned zero. Neither
concerns the artifact a reader downloads: the first says the checkout moved to
the reviewed commit, the second says the process answers on the host. Whether
nginx serves that checkout, whether an \texttt{alias} or a regex location
shadows the file, and whether the bytes on the wire equal the bytes in the commit
are all outside what either check observes.

The replacement gate asserts three things from the public side: that the archive
manifest is served byte-identically to its committed blob; that the timestamp
token over the interoperability record is served byte-identically; and that
every file tracked under the sealed archive directory is either served
byte-identically or refused with the status the access policy requires.

Its first run against production failed:

\begin{lstlisting}
[FAIL] live:archive: 1 of 266 tracked files not delivered as expected
    mismatch docs/tsa/2026-10/verify_cyclonedx_data_loss.py
             (94b28f6d2de183e6 != accd674461443e26)
\end{lstlisting}

The served file was 5{,}025 bytes; the committed one 5{,}170. The archive
directory is marked \texttt{-text} so that a sealed artefact keeps its line
endings; the published mirror under the site tree was written through the
default \texttt{* text=auto eol=lf} rule, which normalised 145 carriage returns
away --- one byte per line. Same text, different bytes.

This is not a cryptographic defect, and calling it one would be the error the
paper is about. It is the same defect as S2 at the level of delivery: the hash in
the manifest no longer described anything downloadable, and every check in the
pipeline was green because every check compared text, or compared the working
tree, which is a rendering of the commit rather than the commit. A failed deploy
is repaired the same hour; a deploy that reports success while serving different
bytes is discovered by a stranger with a hash. A green pipeline is a statement
about the tree; only a request to the wire is a statement about the wire.

\begin{table}[t]
\centering
\small
\begin{tabular}{@{}ll@{}}
\toprule
Status at the commit reviewed for release & Status after the fix \\
\midrule
commit \texttt{6c05152} & commit \texttt{b2910d5} \\
served 5{,}025\,B, 0 CR & served 5{,}170\,B, 145 CR \\
1 of 266 mismatched     & 264 of 266 identical, 2 denied by policy \\
exit 0 (health check only) & exit 0 only after 3 public assertions pass \\
\bottomrule
\end{tabular}
\caption{Delivery state before and after the gate. The two files that are not
byte-identical are \texttt{.sh} artefacts that the access policy denies; for
those the gate requires \emph{refusal}, so a 200 response would fail the deploy.}
\label{tab:deploy}
\end{table}

The repair was a line in \texttt{.gitattributes} and the restoration of one file
from its sealed blob. The manifest was deliberately not regenerated and no token
was re-minted: re-sealing would assert a new fact about a newer file and destroy
the only thing the package proves: on a given date, these were the digests. The mismatch is therefore recorded in a drift note rather than erased.

\subsection{Can the gate fail?}
\label{sec:eval-negative}

A gate that only ever prints PASS is indistinguishable from a script that always
exits zero, so four paths are exercised rather than asserted.

\begin{table}[t]
\centering
\small
\begin{tabular}{@{}lll@{}}
\toprule
Scenario & Expected & Observed \\
\midrule
mock server with files, autoindex, deny list & pass              & \texttt{PASS}, exit 0 \\
same mock, deny rule removed                 & 200 where 403 due & \texttt{policy-regression}, exit 1 \\
production before the byte fix               & one mismatch      & \texttt{1 of 266}, exit 1 \\
repository without the sealed artefacts      & cannot decide     & \texttt{cannot read}, exit 2 \\
\bottomrule
\end{tabular}
\caption{Negative controls for the delivery gate. The third row is the real
defect; the other three exist so that the first row means something.}
\label{tab:gate-controls}
\end{table}

The fourth row mirrors the \texttt{unverifiable} verdict of the record checker.
In both cases the design rule is the same: absence of the inputs required to
decide must produce a distinct, non-success outcome, never a warning attached to
a success.

What the gate proves is deliberately narrow, and its own output says so: for
every file tracked under the sealed archive, the public site returns the
committed blob byte for byte, or refuses it with the status the policy requires.
It does not prove that the content is correct, that the archive is complete with
respect to the world, or that any third party has examined it.

Threats to these measurements are listed with the other limitations in
Section~\ref{sec:limitations}.

% ---------------------------------------------------------------------------
% Provenance notes (not for the paper). Every number above traces to a command;
% if a number changes, the command is re-run rather than the sentence edited.
%
%   tab:verdicts       fibemate-tools: node verifact/bin/verifact.js verify
%                        --config verifact/verifact.json --format json
%                      at fibemate-tools 163efe3 against fibemate b2910d5, both
%                      CLEAN checkouts (git worktree), run 2026-10-30 20:2x:
%                      claimsTotal 2337, codeArtifactClaims 2127,
%                      allJudgedClaims 2207 (= 2127 + 72 dispositioned + 8 denied),
%                      verified 136, unbound 1759, ambiguous 214, refFound 18,
%                      standardRefs 25, sealed 84, ignored 21, denied 8, drift 0,
%                      exit 0; exemptions 19 (6 drift, 13 retained) measured
%                      2026-10-30, base commit a643568.
%                      The earlier figures 2345 / 2135 / 2215 came from a working
%                      tree that carried untracked documents - the defect section
%                      5.3 describes, occurring in this paper's own table. A clean
%                      checkout of the CBOM commit 95056b4 yields 2203 / 2103 /
%                      2112 with drift 15 and exit 1, so the counts are pinned to
%                      b2910d5, the commit whose CBOM figures the paper reports.
%   tab:s0s4           node scripts/claim-evidence/run-e2e.cjs (2026-10-30)
%   tab:s0s4           fibemate-tools/docs/claim-evidence-e2e-2026-10-29.md
%   tab:deploy         6c05152 -> b2910d5; check-live run 2026-10-30:
%                      manifest 30238 B sha256 a6064635…, token 4643 B
%                      sha256 7e3477b0…, archive 264/266 + 2 denied, exit 0
%   tab:gate-controls  scripts/check-live.cjs exit codes: 0 / 1 / 2;
%                      policy-regression path exercised against a mock server
%   145 CR / 5170 B    git cat-file blob HEAD:docs/tsa/2026-10/verify_cyclonedx_data_loss.py
%                      -> 5170 bytes, 145 CR, 145 LF
%   182 KAT / 7-7      fibemate-tauri cargo test --all-features;
%                      fibemate-mobile node --test test/kem-clean-break.node.cjs
% ---------------------------------------------------------------------------
